\chapter{Mathematical Baseline}
\label{ch:mathematics}

\section{Exceptional and extended Lie algebras}

The finite-dimensional exceptional complex simple Lie algebras are
$G_2$, $F_4$, $E_6$, $E_7$, and $E_8$. Their ranks, dimensions, root counts,
and Cartan matrices are established classification data. The affine and
indefinite extensions commonly denoted $E_9$, $E_{10}$, and $E_{11}$ are
Kac--Moody algebras, not additional members of the finite exceptional list
\citep{kac1990infinite}.

For a simply laced finite root system with root lattice $Q$ and weight lattice
$P$, the finite quotient $P/Q$ has order equal to the determinant of the Cartan
matrix. Thus

\begin{equation}
  |P/Q| = \det A,
  \qquad
  \det A_{E_6}=3,
  \quad
  \det A_{E_7}=2,
  \quad
  \det A_{E_8}=1.
\end{equation}

\begin{evidence}{Established result}
These identities classify lattices and representations. They do not, without
an explicit mode map and dynamics, determine transport coefficients, forcing
amplitudes, jet counts, or material response.
\end{evidence}

The $E_8$ lattice is even and unimodular in eight dimensions, and its minimal
vectors form the 240 roots of $E_8$. The lattice also appears in the solution
of the eight-dimensional sphere-packing problem \citep{cohn2017sphere}. None of
these statements supplies a physical coupling merely by reusing the label
$E_8$ in a simulation.

The exceptional Jordan algebra, or Albert algebra, is the 27-dimensional
algebra $J_3(\OO)$ of Hermitian $3\times3$ octonionic matrices. Its exceptional
group ladder depends on which structure is preserved:

\begin{table}[htbp]
  \centering
  \caption{Corrected exceptional-group roles associated with the Albert algebra.}
  \label{tab:albert-group-ladder}
  \begin{tabular}{lll}
    \toprule
    Group & Dimension & Preserved structure \\
    \midrule
    $F_4$ & 52 & Jordan-algebra automorphisms \\
    $E_{6(-26)}$ & 78 & determinant-preserving reduced structure \\
    $E_{7(-25)}$ & 133 & conformal/Freudenthal structure \\
    \bottomrule
  \end{tabular}
\end{table}

Therefore $E_7$ is not the automorphism group of a putative $J_3(S)$ built
from sedenions. The repository audit now evaluates a nonzero octonionic Jordan
product, verifies its identity and commutativity residuals, tests the Jordan
identity, and records the three distinct group roles. This replaces the former
implementation, whose matrix product returned zero for every input.

\section{Cayley--Dickson algebras}

The Cayley--Dickson construction produces real algebras of dimension $2^n$.
The first four are

\begin{equation}
  \RR,
  \qquad
  \CC,
  \qquad
  \HH,
  \qquad
  \OO.
\end{equation}

Algebraic properties are progressively lost under doubling. In particular,
the sedenions contain zero divisors, so the existence of a high-dimensional
Cayley--Dickson algebra does not make it a division algebra or a viable space
of physical states \citep{baez2001octonions,schafer1954composition}.

\begin{evidence}{Established result with a computational illustration}
A finite program can exhibit a zero-divisor pair or test identities on sampled
elements. Such a run illustrates a known theorem; it neither proves the theorem
nor validates a physical use of the algebra.
\end{evidence}

\scriptsize
\begin{longtable}{lrrrrrrrr}
  \caption{Cayley--Dickson property boundary reproduced by the executable audit.}
  \label{tab:cayley-property-boundary}\\
  \toprule
  Algebra & Dim. & Comm. & Assoc. & Alt. & Power & Flex. & Norm & ZD \\
  \midrule
  \endfirsthead
  \toprule
  Algebra & Dim. & Comm. & Assoc. & Alt. & Power & Flex. & Norm & ZD \\
  \midrule
  \endhead
  Real & 1 & Yes & Yes & Yes & Yes & Yes & Yes & No \\
Complex & 2 & Yes & Yes & Yes & Yes & Yes & Yes & No \\
Quaternion & 4 & No & Yes & Yes & Yes & Yes & Yes & No \\
Octonion & 8 & No & No & Yes & Yes & Yes & Yes & No \\
Sedenion & 16 & No & No & No & Yes & Yes & No & Yes \\
Pathion & 32 & No & No & No & Yes & Yes & No & Yes \\
\bottomrule

\end{longtable}
\normalsize

Here ``Power'' means power-associative, ``Flex.'' means flexible, ``Norm''
means multiplicative norm, and ``ZD'' records a nonzero zero divisor. The exact
sedenion witness is $(e_3+e_{10})(e_6-e_{15})=0$.

\section{Modular forms and moonshine}

A holomorphic modular form of integer weight $k$ for $\mathrm{SL}(2,\ZZ)$
is holomorphic on the upper half-plane, satisfies
$f((a\tau+b)/(c\tau+d))=(c\tau+d)^k f(\tau)$, and is holomorphic
at the cusp. For subgroups, every cusp requires its own boundary condition.
The modular invariant $j$ is a modular function: it has a pole at the cusp. The classical modular invariant has
the expansion

\begin{equation}
  j(\tau) = q^{-1} + 744 + 196884q + 21493760q^2 + \cdots,
  \qquad q=e^{2\pi i\tau}.
\end{equation}

The appearance of Monster-group representation dimensions in coefficients of
$j$ led to monstrous moonshine, proved using the Monster
Lie algebra and its twisted denominator identities \citep{conway1979monstrous,borcherds1992monstrous}.

\begin{evidence}{Established result}
Recomputing a few coefficients or special values checks an implementation.
It does not establish that an unrelated physical system has Monster symmetry.
That requires a defined action, observables transforming under it, and tests
that distinguish the symmetry from alternatives.
\end{evidence}

In particular, the theorem identifies graded Monster traces with specified
genus-zero modular functions. It supplies no generic stability operator for a
fluid, fractal, or material evolution equation. The physical claim is rejected
unless a state space, Monster action, evolution operator, invariant or norm,
stability inequality, and matched non-moonshine control are all defined.

\section{Fractal dimensions}

Hausdorff and box-counting dimensions are well-defined mathematical invariants,
but finite-resolution estimators depend on scale selection, sampling, and
regression choices \citep{falconer2004fractal}. A high coefficient of
determination for a log--log fit is not by itself evidence of asymptotic scaling.
The fitted window must be stable under resolution changes and must avoid the
finite-size and discretization plateaus.

\section{The inference boundary}

The established objects in this chapter can parameterize models. They do not
select a model on their own. A defensible cross-domain construction must provide
all of the following:

\begin{enumerate}
  \item a map from the mathematical object to physical degrees of freedom;
  \item equations of motion that use the mapped structure;
  \item units and normalization for every parameter;
  \item a baseline model without the proposed structure;
  \item an observable that differs between the models; and
  \item uncertainty estimates or a convergence study for the comparison.
\end{enumerate}

Without these items, the mathematics may be correct while the physical
interpretation remains an unsupported extrapolation.
